\chapter{Mentor Guide}

This guide is written for the co-advising structure around uafR, water quality, aquaculture/ecosystem context, and machine learning. It is intentionally practical. The goal is to keep students moving while protecting scientific rigor.

\section{Pre-Program Setup Checklist}

\begin{itemize}
  \item Decide whether each student starts with real approved data, teaching data, or both.
  \item Prepare data-use restrictions for unpublished or partner-owned files.
  \item Confirm RStudio and uafR installation path.
  \item Confirm whether live database access is allowed for \code{categorate()} and \code{pubchemProfile()}.
  \item Place any approved GC-MS files in a raw-data folder and record their provenance.
  \item Prepare weekly review times with both co-advising perspectives represented.
  \item Make the dsDNA Core logo and DSU branding assets available if the manual is rendered with final visuals.
\end{itemize}

\section{Daily Check-In Pattern}

A productive daily check-in can be short:
\begin{itemize}
  \item What product is due today?
  \item What file will change?
  \item What evidence is needed before interpretation?
  \item What could go wrong technically?
  \item What claim is currently not supported?
\end{itemize}

\section{Weekly Review Questions}

\begin{longtable}{p{0.15\textwidth}p{0.75\textwidth}}
\toprule
\textbf{Week} & \textbf{Review focus} \\
\midrule
1 & Is the question feasible, WiCCED-relevant, and bounded by available evidence? \\
2 & Can the project folder be understood and rerun from the README? \\
3 & Are variables, units, QA/QC rules, and missingness documented before analysis? \\
4 & Is uafR being used only where chemical evidence supports it, and are claims cautious? \\
5 & Is the ML target valid, is the split defensible, and is there a baseline? \\
6 & Are features frozen, leakage risks addressed, beginner ML understood, and advanced validation/sensitivity outputs interpreted cautiously if used? \\
7 & Do figures and literature synthesis support exactly the claims being made? \\
8 & Can the final package be archived, rerun, and continued by another student? \\
\bottomrule
\end{longtable}

\section{When to Slow Down}

Slow the program if:
\begin{itemize}
  \item The student cannot explain what one row of the dataset represents.
  \item Units, methods, or missing-value codes are unknown.
  \item uafR outputs are being interpreted before validation.
  \item Model performance is discussed before split design.
  \item Figures have claims that are stronger than the data.
  \item The research log is too thin to reconstruct work.
\end{itemize}

\section{When to Accelerate}

Accelerate only after the student has a reproducible pipeline and accurate interpretation habits. Strong students can add:
\begin{itemize}
  \item The advanced ML extension, including site-holdout validation, feature-set sensitivity, permutation importance, and decision-gate review.
  \item A second target variable.
  \item Controlled chemical trait ontology matrices.
  \item Spatial covariates or remote-sensing context, with leakage review.
  \item A reviewer-response appendix.
  \item A short methods supplement suitable for a manuscript or symposium abstract.
\end{itemize}

\section{Supporting Different Backgrounds}

Students with weak R backgrounds should spend more time on Week 2 scripts, folder structure, and console interpretation. Avoid live web enrichment until local scripts run reliably.

Students with limited chemistry background should spend more time on Week 3 and Week 4 evidence types. Require written distinctions among measured variable, tentative annotation, database descriptor, and confirmed identity.

Students with strong data-science backgrounds should be challenged on leakage, sampling design, uncertainty, and scientific claim strength rather than simply adding more complex algorithms.

\section{Real Data, Case Data, or No Data}

The program works in three modes:
\begin{itemize}
  \item \textbf{Real data mode:} Students use approved project data with explicit data-use notes.
  \item \textbf{Case data mode:} Students use de-identified or prepared project-like data.
  \item \textbf{Teaching data mode:} Students use the included synthetic teaching data to learn the workflow while real data are being prepared.
\end{itemize}

Teaching data should be labeled clearly as teaching data. It can support training products, but not environmental claims about Delaware waters.

\section{Protecting Private and Unpublished Data}

Project data may include unpublished sampling locations, partner information, or sensitive infrastructure context. Students should not email raw data, upload private data to public repositories, place unrestricted files in shared folders, or include exact sensitive locations in presentations unless approved. Public-facing outputs should use generalized site labels when required.

\section{Final Package Evaluation}

The final package is acceptable when:
\begin{itemize}
  \item The README states the project question, data restrictions, script order, and expected outputs.
  \item Raw and processed data are separated.
  \item QA/QC decisions are documented.
  \item uafR evidence is validated and source-backed if used.
  \item ML models have baseline comparisons and model cards.
  \item Figures and report claims are consistent.
  \item A clean-session rerun succeeds or blocked steps are documented.
  \item The continuation note identifies next experiments and unresolved risks.
\end{itemize}
